% Scalar SDAC construction and guarantees; stable labels are retained after reorganization.
\chapter{Online Storage Control via Simplex Disturbance-Action Policies}
\label{Section:ChapControl}
\noindent\textbf{Source and author contributions.}
This chapter is based on and extends the preprint
``Online Constrained Control of Storage Systems via Simplex Disturbance-Action
Policies''~\cite{asgari2026sdac}, coauthored with Michael J. Neely.
I formulated the storage-control problem, developed the SDAC policy class
and the original algorithm, and wrote most of the original proofs of the
scalar upper bounds. Neely reviewed the original manuscript and proofs.

The lower-bound construction and proofs in
Section~\ref{ch4:sec:regret-lower-bound} were developed primarily using
\mbox{ChatGPT}, as disclosed in the source manuscript.

\section{Introduction}
\label{ch4:sec:introduction}

Chapter~\ref{Section:ChapEnergyharvest} studies online allocation with
stochastic energy arrivals in $[0,1]$ and a fixed-action benchmark. Here,
arrivals may follow any sequence in $[0,1]$, each action is selected before
the current arrival is observed, and convex costs may depend jointly on
storage and allocation. Each fixed comparator policy is evaluated on its
own storage trajectory. Its actions may vary with the arrival history.

The scalar model has storage decay and enough capacity to avoid saturation.
A \emph{Simplex Disturbance-Action Control} (SDAC) policy assigns
nonnegative fractions of each arrival to withdrawal delays, with total
fraction at most one. For a fixed SDAC policy, the storage level and allocation
are affine functions of its allocation fractions.
Lemmas~\ref{ch4:lem:sdac-state-decomp}--\ref{ch4:lem:sdac-surrogate} show that
the policy never withdraws more resource than is available and
that its cost is convex in these policy parameters.

The \emph{entropic SDAC controller} uses entropic mirror descent and clips
requested allocations to the resource currently available. At each round,
the surrogate cost is evaluated at the state and action that would result
if the current allocation fractions had been used from the start with the
same arrival history. We bound regret using these hypothetical costs, then
bound their difference from the costs actually incurred as the fractions
change and allocations are clipped.
After each action, the controller uses the revealed full cost function to
obtain a subgradient at this hypothetical state and action.

\subsection{Contributions and scope}

The entropic SDAC controller is feasible on every round. For a chosen memory
length $H$, setting the learning step size according to the horizon $T$ gives
$O(\sqrt{T\log(H+1)})$ regret against policies with memory $H$, with the
retention coefficient and cost bounds held fixed. We approximate
each infinite-memory policy by discarding its withdrawal weights for delays
greater than $H$. Storage decay gives a bound on the resulting cost difference
that decreases geometrically with $H$. Choosing $H$ proportional to $\log T$
as specified in~\eqref{ch4:eq:def-infinite-memory-length} gives
$O(\sqrt{T\log\log T})$ regret against the best fixed infinite-memory SDAC
policy. This class includes policies that withdraw a constant fraction of
the currently stored resource each round: $u_t=kx_t$, where $0\le k\le\alpha$.
Each such policy can be represented as an infinite-memory SDAC policy.

For $G_\alpha\coloneqq\alpha(L_x/(1-\alpha)+L_u)$ and horizons satisfying
\eqref{ch4:eq:sdac-lower-horizon},
Theorem~\ref{ch4:thm:sdac-lower-bound} gives a worst-case expected
regret lower bound of $G_\alpha\sqrt{T}/8$ for every causal feasible
controller, including randomized controllers.
For fixed $H,\alpha,L_x,L_u$, this matches the upper bound's
$\sqrt{T}$ dependence on the horizon $T$. For fixed $L_x,L_u>0$, as
$\alpha\uparrow1$ and with the same horizon condition, the lower bound is
$\Omega((1-\alpha)^{-1}\sqrt T)$.
The upper and lower bounds still differ by factors of
$\sqrt{\log(H+1)}$ in memory dependence and $(1-\alpha)^{-1/2}$ in
retention dependence. Closing these gaps remains open.

For energy-harvesting storage, the state, arrivals, and actions represent
stored, harvested, and allocated energy, respectively. The model also
applies to other resources with proportional storage loss.

The joint manuscript~\cite{asgari2026sdac} introduces the scalar SDAC policy
class and entropic controller. It also gives the equivalence
between clipping and projection, the finite- and infinite-memory upper
bounds, the fixed-fraction inclusion, and the regret lower bound.

Sections~\ref{ch4:sec:scalar} and~\ref{ch4:sec:scalar-performance}
give the controller and its finite- and infinite-memory guarantees.
Section~\ref{ch4:sec:regret-lower-bound} gives the lower bound and compares
its parameter dependence with the upper bound.
Section~\ref{ch4:subsec:scalar-entropy} proves the projection interpretation of action clipping.

\subsection{Related work}
\label{ch4:sec:related-work}

\paragraph{Online control}
Goulart, Kerrigan, and Maciejowski
\cite{goulart2006optimization} express control actions as affine functions
of past disturbances. For convex costs, this makes finite-horizon policy
optimization convex while enforcing constraints involving both the state
and the action for every allowed bounded disturbance sequence. Agarwal
et al.\ combine a stabilizing linear controller with learned responses to past disturbances
and obtain $\widetilde O(\sqrt{T})$ regret against strongly stable linear
controllers \cite{agarwal2019online}. Golowich et al.\
\cite{golowich2024population} adapt policies based on past disturbances to
population dynamics, using nonnegative policy parameters subject to
normalization constraints. They use entropy-based policy updates and bound
the cost difference caused by changing the policy over time. Their state
records population proportions, and the available actions do not depend
on the current state. Here, the simplex weights distribute arrivals across
withdrawal times, and feasibility requires
$0\le u_t\le\alpha x_t$.

Li et al.\ consider separate affine constraints on the state and action
and assume a linear policy that always maintains a fixed positive margin
from every constraint boundary \cite{li2021onlineConstrained}. They
project onto a restricted set of policy parameters that preserves
feasibility. Such a margin need not exist here: when storage is empty,
the only feasible allocation is zero.

Liu et al.\ \cite{liu2023constrainedControl} allow convex constraints
involving both the state and the action. They bound accumulated violation
for changing constraints and violation at each round for fixed constraints.
Projection can satisfy fixed affine constraints exactly
\cite[Table~1]{liu2023constrainedControl}. Our constraint is fixed and
affine. Fixed SDAC policies and the clipping
rule~\eqref{ch4:eq:alg-feasible-action} satisfy it exactly.

When costs depend on past decisions, the analysis can first consider the
cost of repeatedly using the current decision, then bound the additional
cost caused by changing decisions over time
\cite{anava2015memory,kumar2023unbounded}. Our tracking estimate bounds
this additional cost using nonnegative arrivals and decay in the storage
dynamics, while accounting for allocations clipped to the available resource.
Kumar, Dean, and Kleinberg \cite{kumar2023unbounded} also give a lower
bound for costs whose dependence on past decisions decreases geometrically,
using a Euclidean norm on decision histories. For sufficiently long
horizons and discount factors approaching one, their bound scales as
$\sqrt T$ divided by one minus the discount factor. This parallels the
retention dependence of Theorem~\ref{ch4:thm:sdac-lower-bound}
under \eqref{ch4:eq:sdac-lower-horizon}. The storage result applies to
causal feasible controllers under the state-dependent action constraint
and compares with fixed SDAC policies on their own trajectories.

\paragraph{Energy harvesting}
Under i.i.d.\ arrivals, Shaviv and \"Ozg\"ur
\cite{shaviv2016universally} allocate a fixed fraction of stored energy
each round and keep the long-run expected transmission rate within a
constant of the optimum. Arafa et al.\ \cite{arafa2018online} extend such
policies to increasing, concave utility functions under i.i.d.\ arrivals.
For Bernoulli arrivals that either fully recharge the battery or add no
energy, they characterize allocations that decrease or stay constant
between recharges. Zibaeenejad and Chen \cite{zibaeenejad2019lookahead}
derive an optimal policy under Bernoulli arrivals with advance knowledge
of a fixed number of future arrivals. The policy uses a decreasing
allocation sequence when none of those arrivals recharges the battery.
These policies and SDAC both distribute arrivals across future allocations.
The exact inclusion in
Corollary~\ref{ch4:cor:fraction-as-sdac} concerns feasible fixed-fraction
policies under the chapter's own dynamics.

Yu and Neely \cite{Yu2019} study utility maximization when current
operating conditions are revealed only after the allocation decision.
Their journal version allows conditions other than energy arrivals to be
correlated over time and to have changing distributions, while energy
arrivals are independent and have a common mean
\cite[Sec.~V, Thm.~4]{Yu2019}. Asgari and Neely \cite{10.1145/3428337}
report expected regret $O(\sqrt{T})$ with battery capacity $O(\sqrt{T})$
for fixed dimension, with i.i.d.\ energy arrivals and time-varying convex
losses under their stated independence assumptions.
Chapter~\ref{Section:ChapEnergyharvest} makes the conditional-mean
condition~\eqref{eq:energy:conditional-mean} explicit for history-dependent
costs. Both papers compare expected performance against
constant allocations whose total energy use per round does not exceed the
mean energy arrival. These comparison allocations need not satisfy the
battery constraint on every realized arrival sequence. Neither objective
depends directly on the battery state.
Lin et al.\ \cite{lin2025optimal} allow
arbitrary energy arrivals and changing objectives that measure the freshness
of received information. They observe the current arrival and
reward function before allocation. They bound the worst-case ratio between
the optimal reward with the entire input sequence known in advance and
their controller's reward, under the same battery-capacity and
allocation-rate constraints.
Our setting combines arbitrary bounded arrivals, decisions made before the
current arrival and cost are revealed, and convex costs depending jointly
on storage and allocation. We measure additive regret against fixed SDAC
policies evaluated on their own trajectories.

\paragraph{Storage management}
Qin et al.\ \cite{qin2016online} study storage control under uncertain
prices and net demand. Their controller balances current operating cost
against a term depending on stored energy. Under their assumptions on the
input process, it respects the storage constraints and bounds how far its
long-run expected average cost exceeds the optimal value. Kim et al.\
\cite{kim2016storage} study energy routing in networks with storage that
loses energy over time. They give competitive guarantees for the value of
meeting demand minus production cost. Their guarantee including production
costs assumes shared storage and advance knowledge of demands. These
papers bound long-run average cost or ratios of total performance. This
chapter bounds the cumulative cost difference relative to the best fixed
policy chosen in hindsight.

\paragraph{Queueing and network optimization}
Classical queueing control studies throughput and queue stability through dynamic
service and scheduling decisions
\cite{tassiulas1990stability,tassiulas1993dynamic,mckeown1999achieving}, while
stochastic network optimization studies tradeoffs between time-average performance
and average queue backlog \cite{neely2010stochastic}.

Khojastepour and Sabharwal \cite{khojastepour2004delay} study schedulers
that split each arrival among service times within a maximum delay. For
their class with fixed allocation fractions, the service rate is a fixed
weighted sum of current and past arrivals. The fractions are nonnegative
and sum to one. SDAC permits their sum to be at most one. Under i.i.d.\
arrivals, spreading each arrival equally across the allowed service slots
minimizes long-run average transmit power over a Gaussian channel within
this class of policies with fixed allocation fractions. Their objective
uses a fixed cost function.

\paragraph{Inventory and reservoir management}
Yang et al.\ \cite{yang2020online} study purchasing to meet demand under
inventory constraints, observing the current demand and price before
acting. They obtain the best possible worst-case ratios between their
cost and the minimum cost with all future demands and prices known.
Hihat et al.\ \cite{hihat2023inventory} study inventory systems where
unsold stock carries over between periods and may expire. They establish
feasibility and $O(\sqrt{T})$ expected regret relative to the best fixed
replenishment target. Their guarantee requires a uniformly positive
conditional probability, given the past, that demand for every product is
at least a fixed positive amount. Hihat and Fermanian
\cite{hihat2024online} subsequently develop a method for learning
replenishment targets for more general inventory dynamics, without a
regret guarantee for that method. Classical reservoir models study how
random inflows and prescribed release rules determine the long-run
distribution of stored water \cite{moran1956continuous}.

\section{Model, policies, and algorithm}
\label{ch4:sec:scalar}

We follow the shared conventions in Section~\ref{sec:intro:notation}.

\subsection{Model}
\label{ch4:sec:model-setup}

Fix a positive integer $T$ and a known retention coefficient
$\alpha\in(0,1)$. The state $x_t\in\Real_+$ is the resource level at the
beginning of time $t$. The system starts empty:
\begin{equation}
\label{ch4:eq:model-init}
x_1=0.
\end{equation}
At each time $t\in[T]$, the controller chooses an action
$u_t\in\Real_+$ satisfying
\begin{equation}
\label{ch4:eq:model-feasibility}
0\le u_t\le\alpha x_t.
\end{equation}
The system then receives an arrival $w_t\in[0,1]$ and evolves according to
\begin{equation}
\label{ch4:eq:model-dynamics}
x_{t+1}=\alpha x_t-u_t+w_t.
\end{equation}
Here, $1-\alpha$ is the fraction lost per period, and
\eqref{ch4:eq:model-feasibility} limits allocation to the resource remaining
after decay.

At time $t$, the controller chooses $u_t$ before observing $w_t$ or $c_t$. After the
action is selected, the arrival and the entire cost function are revealed, the
controller incurs $c_t(x_t,u_t)$, and the state is updated according to
\eqref{ch4:eq:model-dynamics}.

For every feasible trajectory, \eqref{ch4:eq:model-dynamics},
\eqref{ch4:eq:model-feasibility}, and $w_t\in[0,1]$ give
$0\le x_{t+1}\le\alpha x_t+1$. Iterating from
\eqref{ch4:eq:model-init} gives
\begin{equation}
\label{ch4:eq:model-state-bound}
0\le x_t
\le\sum_{j=0}^{t-2}\alpha^j
=\frac{1-\alpha^{t-1}}{1-\alpha}
\le\frac{1}{1-\alpha},
\qquad t\in[T+1].
\end{equation}
The sum is empty at $t=1$.
Together with \eqref{ch4:eq:model-feasibility}, this bound places every feasible
state--action pair in the compact set
\begin{equation}
\label{ch4:eq:model-domain}
\mathcal{D}_\alpha
\coloneqq
\left\{
(x,u):
0\le x\le\frac{1}{1-\alpha},
0\le u\le\alpha x
\right\}.
\end{equation}

\begin{Thm:Remark}[Storage capacity]
\label{ch4:rem:storage-capacity}
A physical storage capacity $B\ge1/(1-\alpha)$ suffices to avoid overflow
on every feasible trajectory starting at zero. Thus the same results hold for
$x_{t+1}=\min\{B,\alpha x_t-u_t+w_t\}$ at these capacities. With smaller
storage capacities, overflow can make the storage level a nonaffine function
of the policy fractions. A convex cost in storage and allocation then need
not remain convex in those fractions. The learning analysis below relies on
that convexity.
\end{Thm:Remark}

Each revealed cost function $c_t:\Real_+^2\to\mathbb{R}$ is convex in
$(x,u)$ and satisfies the following Lipschitz condition on $\mathcal{D}_\alpha$:
there are constants $L_x,L_u>0$ such that
\begin{equation}
\label{ch4:eq:model-lipschitz}
|c_t(x,u)-c_t(x',u')|
\le
L_x|x-x'|+L_u|u-u'|
\end{equation}
for all $(x,u),(x',u')\in\mathcal{D}_\alpha$.
At every $(x,u)\in\mathcal{D}_\alpha$, we also assume access to a cost subgradient
\begin{equation}
\label{ch4:eq:model-cost-subgradient}
\boldsymbol{\zeta}_t(x,u)
\coloneqq
\bigl(\zeta_t^x(x,u),\zeta_t^u(x,u)\bigr)^\top
\in\partial c_t(x,u)
\end{equation}
satisfying
\begin{equation}
\label{ch4:eq:model-subgradient-bounds}
|\zeta_t^x(x,u)|\le L_x,
\qquad
|\zeta_t^u(x,u)|\le L_u.
\end{equation}
These subgradient bounds hold, for example, when $c_t$ has a convex
extension to an open neighborhood of $\mathcal{D}_\alpha$ satisfying
\eqref{ch4:eq:model-lipschitz}. See
\cite[Thm.~9.13]{rockafellar1998variational}.

We use $x_t$ and $u_t$ without parameter arguments for the trajectory generated by
the online controller. Its cumulative cost is
\begin{equation}
\label{ch4:eq:def-cumulative-cost}
J_T^{\mathrm{on}}
\coloneqq
\sum_{t=1}^T c_t(x_t,u_t).
\end{equation}


Our objective is to satisfy \eqref{ch4:eq:model-feasibility} at every time and achieve
sublinear regret against the fixed policies defined next.

\subsection{Finite-memory SDAC policies and regret benchmark}
\label{ch4:sec:sdac}

DAC uses a finite linear filter of past disturbances
\cite{agarwal2019online}, building on affine disturbance-feedback
parameterizations \cite{goulart2006optimization}. In our scalar storage
model, this filter has the following form:
\begin{equation}
\label{ch4:eq:dac-param}
u_t \coloneqq \sum_{i=1}^{H} \theta_i w_{t-i},
\end{equation}
where $H$ is a positive integer (the memory length) and
$\boldsymbol{\theta}\coloneqq(\theta_1,\ldots,\theta_H)^\top\in\mathbb{R}^H$ is the vector
of filter coefficients.

An unrestricted filter can request a negative allocation or more than
$\alpha x_t$. SDAC restricts the coefficients so that each fixed policy is
feasible and its surrogate costs are convex.


Define the subprobability simplex
\begin{equation}
\label{ch4:eq:def-simplex}
\mathcal{M}_H
\coloneqq
\left\{
\mathbf{m}\in\Real_+^H:
\|\mathbf{m}\|_1\le1
\right\}.
\end{equation}
With the convention $w_\tau\coloneqq0$ for $\tau\le0$, each
$\mathbf{m}\in\mathcal{M}_H$ defines an SDAC policy through
\begin{equation}
\label{ch4:eq:def-sdac-policy}
u_t(\mathbf{m})
\;\coloneqq\;
\sum_{i=1}^{H} m_i\alpha^i w_{t-i}
=
\inner{\boldsymbol{\phi}_t^u}{\mathbf{m}},
\end{equation}
where
\begin{equation}
\label{ch4:eq:def-phi-u}
\boldsymbol{\phi}_t^u
\coloneqq
\bigl(\alpha w_{t-1},\alpha^2 w_{t-2},\ldots,\alpha^H w_{t-H}\bigr)^\top
\in\mathbb{R}^H.
\end{equation}
Here, $m_i\alpha^i w_s$ is the amount from arrival $w_s$ used $i$ periods later.

Unless stated otherwise, SDAC policies have finite memory, with parameter
domain $\mathcal M_H$.
The \emph{fixed-policy trajectory}
$\{(x_t(\mathbf m),u_t(\mathbf m))\}_{t=1}^T$ is the unique trajectory
obtained by using $\mathbf m$ at every time.
Its cumulative cost is
\begin{equation}
\label{ch4:eq:def-fixed-policy-cost}
J_T(\mathbf{m})
\coloneqq
\sum_{t=1}^T c_t\bigl(x_t(\mathbf{m}),u_t(\mathbf{m})\bigr).
\end{equation}
For a fixed positive integer $H$, the regret against the best fixed
SDAC policy is
\begin{equation}
\label{ch4:eq:def-regret}
\Reg_T(\mathcal M_H)
\coloneqq
J_T^{\mathrm{on}}
-\min_{\mathbf{m}\in\mathcal{M}_H}J_T(\mathbf{m}).
\end{equation}
The comparison uses the same realized arrival and cost sequence for the
online controller and all fixed policies.

\begin{Thm:Lemma}[State decomposition under SDAC]
\label{ch4:lem:sdac-state-decomp}
Fix a positive integer $H$.
For $t\in[T]$ and $i\in[H]$, define
\begin{equation}
\label{ch4:eq:sdac-residual-definitions}
r_t(i)
\coloneqq
\sum_{j=1}^{i}\alpha^{j-1}w_{t-j},
\qquad
r_t
\coloneqq
\sum_{j=1}^{t-1}\alpha^{j-1}w_{t-j},
\end{equation}
where the empty sum defining $r_1$ is zero. Then, for every
$\mathbf{m}\in\mathcal{M}_H$ and $t\in[T]$, the fixed-policy state associated
with $\mathbf{m}$ satisfies
\begin{equation}
\label{ch4:eq:sdac-state-decomp}
x_t(\mathbf{m})
=
\sum_{i=1}^H m_i r_t(i)
+\left(1-\sum_{i=1}^H m_i\right)r_t.
\end{equation}
Equivalently, with
\begin{equation}
\label{ch4:eq:def-phi-x}
\boldsymbol{\phi}_t^x
\coloneqq
\bigl(r_t(1)-r_t,\ldots,r_t(H)-r_t\bigr)^\top,
\end{equation}
the state and action satisfy
\begin{equation}
\label{ch4:eq:surr-affine-pair}
x_t(\mathbf{m})=r_t+\inner{\boldsymbol{\phi}_t^x}{\mathbf{m}},
\qquad
u_t(\mathbf{m})=\inner{\boldsymbol{\phi}_t^u}{\mathbf{m}}.
\end{equation}
Thus, both are affine in $\mathbf m$.
\end{Thm:Lemma}

\begin{proof}
Shifting the index in \eqref{ch4:eq:sdac-residual-definitions} gives,
for $1\le t<T$ and $i\in[H]$,
\begin{align*}
r_{t+1}&=w_t+\alpha r_t, \\
r_{t+1}(i)
&=w_t+\sum_{j=1}^{i-1}\alpha^j w_{t-j}
=w_t+\alpha r_t(i)-\alpha^i w_{t-i}.
\end{align*}
With empty sums equal to zero and $w_\tau=0$ for $\tau\le0$, the
recurrences also hold at $t=1$ and $i=1$. Since the system starts empty,
$x_1(\mathbf m)=r_1(i)=r_1=0$.
If \eqref{ch4:eq:sdac-state-decomp} holds at $t$, then
\eqref{ch4:eq:model-dynamics} and \eqref{ch4:eq:def-sdac-policy} give
\begin{align*}
x_{t+1}(\mathbf{m})
&=\alpha x_t(\mathbf{m})-u_t(\mathbf{m})+w_t\\
&=\sum_{i=1}^H m_i
\bigl(\alpha r_t(i)-\alpha^i w_{t-i}\bigr)
+\left(1-\sum_{i=1}^H m_i\right)\alpha r_t+w_t\\
&=\sum_{i=1}^H m_i\bigl(r_{t+1}(i)-w_t\bigr)
+\left(1-\sum_{i=1}^H m_i\right)(r_{t+1}-w_t)+w_t\\
&=\sum_{i=1}^H m_i r_{t+1}(i)
+\left(1-\sum_{i=1}^H m_i\right)r_{t+1}.
\end{align*}
This completes the induction.
\end{proof}



\begin{Thm:Lemma}[Feasibility of SDAC policies]
\label{ch4:lem:sdac-feasibility}
For every $\mathbf{m}\in\mathcal{M}_H$, the fixed SDAC policy
indexed by $\mathbf{m}$ is feasible:
\[
0\le u_t(\mathbf{m})\le \alpha x_t(\mathbf{m})
\qquad \text{for every }t\in[T].
\]
\end{Thm:Lemma}

\begin{proof}
By \eqref{ch4:eq:sdac-residual-definitions}, for $t\in[T]$ and $i\in[H]$,
\begin{align*}
 \alpha r_t(i)-\alpha^iw_{t-i}
 &=\sum_{j=1}^{i-1}\alpha^jw_{t-j}\ge0,
 \qquad
 r_t=\sum_{j=1}^{t-1}\alpha^{j-1}w_{t-j}\ge0.
\end{align*}
Substituting \eqref{ch4:eq:sdac-state-decomp} and
\eqref{ch4:eq:def-sdac-policy},
\begin{align*}
u_t(\mathbf m)
 &=\sum_{i=1}^Hm_i\alpha^iw_{t-i}\ge0,\\
\alpha x_t(\mathbf m)-u_t(\mathbf m)
 &=\sum_{i=1}^Hm_i\bigl(\alpha r_t(i)-\alpha^iw_{t-i}\bigr)
   +\alpha\left(1-\sum_{i=1}^Hm_i\right)r_t
\overset{\mathrm{(A)}}{\ge}0.
\end{align*}
Here (A) uses the nonnegative residuals above and the nonnegative weights
$m_i$ and $1-\sum_i m_i$ from \eqref{ch4:eq:def-simplex}.
\end{proof}



\begin{Thm:Lemma}[Convexity and Lipschitz continuity of the surrogate cost]
\label{ch4:lem:sdac-surrogate}
For $t\in[T]$, define the fixed-policy surrogate cost
\begin{equation}
\label{ch4:eq:def-surrogate-cost}
\ell_t(\mathbf{m})
\coloneqq
c_t\bigl(x_t(\mathbf{m}),u_t(\mathbf{m})\bigr),
\qquad
\mathbf{m}\in\mathcal{M}_H,
\end{equation}
and
\begin{equation}
\label{ch4:eq:def-surrogate-G}
G_\alpha
\coloneqq
\alpha\left(\frac{L_x}{1-\alpha}+L_u\right).
\end{equation}
Then $\ell_t$ is convex on $\mathcal{M}_H$ and, for all
$\mathbf{m},\mathbf{m}'\in\mathcal{M}_H$,
\[
|\ell_t(\mathbf{m})-\ell_t(\mathbf{m}')|
\le
G_\alpha\|\mathbf{m}-\mathbf{m}'\|_1.
\]
For $\mathbf{m}\in\mathcal{M}_H$, use \eqref{ch4:eq:model-cost-subgradient} to define
\begin{equation}
\label{ch4:eq:surr-subgradient}
\begin{aligned}
(\zeta_t^x,\zeta_t^u)^\top
&\coloneqq
\boldsymbol{\zeta}_t\bigl(x_t(\mathbf{m}),u_t(\mathbf{m})\bigr), \\
\mathbf{g}_t
&\coloneqq
\zeta_t^x\boldsymbol{\phi}_t^x
+\zeta_t^u\boldsymbol{\phi}_t^u.
\end{aligned}
\end{equation}
Then
\[
\mathbf{g}_t\in\partial\ell_t(\mathbf{m}),
\qquad
\|\mathbf{g}_t\|_\infty\le G_\alpha.
\]
\end{Thm:Lemma}

\begin{proof}
By \eqref{ch4:eq:sdac-residual-definitions} and $w_\tau=0$ for $\tau\le0$,
we have $0\le r_t(i)\le r_t$, with $r_t(i)=r_t$ when $i\ge t-1$.
Thus \eqref{ch4:eq:def-phi-x} gives, for $i\in[H]$,
\[
\begin{aligned}
|\phi_{t,i}^x|
 &=|r_t(i)-r_t|=r_t-r_t(i)
   =\sum_{j=i+1}^{t-1}\alpha^{j-1}w_{t-j}\\
 &\le\sum_{j=i+1}^{\infty}\alpha^{j-1}
   =\frac{\alpha^i}{1-\alpha}
   \le\frac{\alpha}{1-\alpha},\\
0\le\phi_{t,i}^u&=\alpha^i w_{t-i}\le\alpha^i\le\alpha.
\end{aligned}
\]
Here $w_s\in[0,1]$, $0<\alpha<1$, and the finite sum is empty
when $i\ge t-1$. Hence
\begin{equation}
\label{ch4:eq:surr-coefficient-norms}
\|\boldsymbol{\phi}_t^x\|_\infty\le\frac{\alpha}{1-\alpha},
\qquad
\|\boldsymbol{\phi}_t^u\|_\infty\le\alpha.
\end{equation}
Lemma~\ref{ch4:lem:sdac-feasibility} and
\eqref{ch4:eq:model-state-bound} place the fixed-policy pairs in
$\mathcal D_\alpha$, where the cost and subgradient bounds apply.
The affine formula \eqref{ch4:eq:surr-affine-pair} and
Lemma~\ref{lem:intro:convex-comparison} give convexity. For
$\mathbf m,\mathbf m'\in\mathcal M_H$,
\[
\begin{aligned}
\ell_t(\mathbf m')-\ell_t(\mathbf m)
&\overset{\mathrm{(A)}}{\ge}
 \zeta_t^x[x_t(\mathbf m')-x_t(\mathbf m)]
 +\zeta_t^u[u_t(\mathbf m')-u_t(\mathbf m)]\\
&\overset{\mathrm{(B)}}{=}
 \zeta_t^x\inner{\boldsymbol\phi_t^x}{\mathbf m'-\mathbf m}
 +\zeta_t^u\inner{\boldsymbol\phi_t^u}{\mathbf m'-\mathbf m}\\
&=\inner{\zeta_t^x\boldsymbol\phi_t^x+\zeta_t^u\boldsymbol\phi_t^u}
        {\mathbf m'-\mathbf m}\\
&\overset{\mathrm{(C)}}{=}\inner{\mathbf g_t}{\mathbf m'-\mathbf m}.
\end{aligned}
\]
Here (A) uses the subgradient in \eqref{ch4:eq:model-cost-subgradient};
(B) substitutes \eqref{ch4:eq:surr-affine-pair};
(C) is \eqref{ch4:eq:surr-subgradient}.
Thus $\mathbf g_t\in\partial\ell_t(\mathbf m)$, as in
\eqref{eq:intro:affine-subgradient}, and
\[
\begin{aligned}
\|\mathbf g_t\|_\infty
&=\|\zeta_t^x\boldsymbol\phi_t^x+\zeta_t^u\boldsymbol\phi_t^u\|_\infty\\
&\le |\zeta_t^x|\,\|\boldsymbol\phi_t^x\|_\infty
     +|\zeta_t^u|\,\|\boldsymbol\phi_t^u\|_\infty\\
&\overset{\mathrm{(A)}}{\le}
 L_x\|\boldsymbol\phi_t^x\|_\infty+L_u\|\boldsymbol\phi_t^u\|_\infty\\
&\overset{\mathrm{(B)}}{\le}L_x\frac{\alpha}{1-\alpha}+L_u\alpha
 =G_\alpha.
\end{aligned}
\]
Here (A) uses \eqref{ch4:eq:model-subgradient-bounds}, (B) uses
\eqref{ch4:eq:surr-coefficient-norms}, and the equality is
\eqref{ch4:eq:def-surrogate-G}. Therefore
\[
\begin{aligned}
\ell_t(\mathbf m)-\ell_t(\mathbf m')
&\le\inner{\mathbf g_t}{\mathbf m-\mathbf m'}\\
&\le\|\mathbf g_t\|_\infty\|\mathbf m-\mathbf m'\|_1\\
&\le G_\alpha\|\mathbf m-\mathbf m'\|_1.
\end{aligned}
\]
Interchanging $\mathbf m$ and $\mathbf m'$ gives the Lipschitz bound,
as in Lemma~\ref{lem:intro:lipschitz}.
\end{proof}

\subsection{Entropic SDAC controller}
\label{ch4:sec:online-algorithm}

The controller updates the allocation fractions
$\mathbf m_t\in\mathcal M_H=S_H$ using entropic online mirror descent (OMD)
\cite{BECK2003167} with step size $\eta>0$, the unnormalized negentropy
$\Psi$, and generalized KL divergence $D_\Psi$ from
\eqref{eq:intro:negentropy}--\eqref{eq:intro:generalized-kl}.
The memory $H$ and step size $\eta$ are supplied parameters.
The divergence is defined for $\mathbf m\in\mathcal M_H$ and
$\mathbf q>\mathbf0$. Initialize
\begin{equation}
\label{ch4:eq:alg-init}
\mathbf{m}_1\coloneqq\frac{1}{H+1}\mathbf{1}.
\end{equation}

\paragraph{Action selection}
Fixed-policy feasibility need not hold when $\mathbf{m}_t$ varies: the
requested action $\inner{\boldsymbol{\phi}_t^u}{\mathbf{m}_t}$ can exceed
$\alpha x_t$. For $H=2$, $w_1=1$, and $w_2=0$, take
$\mathbf m_2\coloneqq(2/3,1/3)^\top$ and $\mathbf m_3\coloneqq(1/3,2/3)^\top$ in
$\mathcal M_2$. Then
\[
 x_2=1,\qquad u_2=\frac{2\alpha}{3},\qquad x_3=\frac{\alpha}{3},
 \qquad
 \inner{\boldsymbol\phi_3^u}{\mathbf m_3}
 =\frac{2\alpha^2}{3}>\alpha x_3=\frac{\alpha^2}{3}.
\]
The controller therefore clips the requested action:
\begin{equation}
\label{ch4:eq:alg-feasible-action}
u_t
\coloneqq
\min\!\left\{
\inner{\boldsymbol{\phi}_t^u}{\mathbf{m}_t},
\alpha x_t
\right\}.
\end{equation}
By \eqref{ch4:eq:def-phi-u}, $\phi_{t,i}^u=\alpha^i w_{t-i}\ge0$ for every
$i\in[H]$, because $\alpha>0$ and the arrivals are nonnegative. Thus
\eqref{ch4:eq:alg-feasible-action} and $\mathbf m_t\in\mathcal M_H$ give,
whenever $x_t\ge0$,
\[
0\le u_t\le\alpha x_t,
\qquad x_{t+1}=\alpha x_t-u_t+w_t\ge w_t\ge0.
\]
Since $x_1=0$, the online trajectory is feasible. The learning update uses
the surrogate subgradient at $\mathbf m_t$. Clipping only changes the
executed action. Lemma~\ref{ch4:lem:availability-projection} gives another way
to obtain the clipped allocation: among admissible SDAC fractions whose
requested allocation does not exceed $\alpha x_t$, choose those minimizing
the generalized KL divergence from the current fractions. The resulting
allocation equals the clipped request.
\paragraph{Surrogate cost and OMD update}
After $w_t$ and $c_t$ are revealed, the controller evaluates the surrogate
subgradient $\mathbf g_t=\zeta_t^x\boldsymbol\phi_t^x+\zeta_t^u\boldsymbol\phi_t^u$
from \eqref{ch4:eq:surr-subgradient}, with the cost subgradient evaluated at
$(x_t(\mathbf m_t),u_t(\mathbf m_t))$, and sets
\begin{equation}
\label{ch4:eq:alg-omd-argmin}
\mathbf{m}_{t+1}
\coloneqq
\argmin_{\mathbf{m}\in\mathcal{M}_H}
\left\{
\eta\inner{\mathbf{g}_t}{\mathbf{m}}
+D_\Psi(\mathbf{m},\mathbf{m}_t)
\right\}.
\end{equation}
Apply Lemma~\ref{ch4:lem:omd-closed-form} with $d=H$, $Q=S_H=\mathcal M_H$,
$q=\mathbf m_t$, and $g=\mathbf g_t$. Its conditions are $\eta>0$ and
$\mathbf m_t>\mathbf0$. The latter holds initially by
\eqref{ch4:eq:alg-init} and is preserved by the following two-step update:
\begin{equation}
\label{ch4:eq:alg-omd-mult}
\widetilde{m}_{t+1,i}
\coloneqq
m_{t,i}\exp\!\bigl(-\eta g_{t,i}\bigr),
\qquad
i\in[H].
\end{equation}
Then project onto $\mathcal M_H$:
\begin{equation}
\label{ch4:eq:alg-omd-proj}
\mathbf{m}_{t+1}
=\argmin_{\mathbf{m}\in\mathcal{M}_H}
D_\Psi(\mathbf{m},\widetilde{\mathbf{m}}_{t+1})
=\frac{\widetilde{\mathbf{m}}_{t+1}}
{\max\{1,\|\widetilde{\mathbf{m}}_{t+1}\|_1\}}.
\end{equation}


\section{Performance guarantees}
\label{ch4:sec:scalar-performance}

\subsection{Finite-memory regret guarantee}
\label{ch4:sec:regret-analysis}

For every $\mathbf m\in\mathcal M_H$, the fixed-policy cost
\eqref{ch4:eq:def-fixed-policy-cost} and surrogate
\eqref{ch4:eq:def-surrogate-cost} give the decomposition
\begin{equation}
\label{ch4:eq:regret-decomposition}
\begin{aligned}
 J_T^{\mathrm{on}}-J_T(\mathbf m)
 &=\sum_{t=1}^T\bigl[c_t(x_t,u_t)-\ell_t(\mathbf m_t)\bigr]\\
 &\quad+\sum_{t=1}^T\bigl[\ell_t(\mathbf m_t)-\ell_t(\mathbf m)\bigr].
\end{aligned}
\end{equation}
The first sum, called the tracking difference, compares the controller's
actual cost with the cost it would incur at each round if its current
allocation fractions had been used from the start with the same arrival
history. The second sum, called surrogate regret, compares these hypothetical
costs with those of the fixed policy $\mathbf m$.
Bounding these terms gives the following guarantee.

\begin{Thm:Theorem}[Regret against fixed SDAC policies]
\label{ch4:thm:regret-bound}
Fix $\alpha\in(0,1)$ and positive integers $H,T$.
Assume the model and feedback of Subsection~\ref{ch4:sec:model-setup},
with $x_1=0$, $w_t\in[0,1]$, and convex costs satisfying the Lipschitz
and accessible subgradient bounds with constants $L_x,L_u>0$.
For any step size $\eta>0$, the controller in
Subsection~\ref{ch4:sec:online-algorithm}, initialized by \eqref{ch4:eq:alg-init}, satisfies
\begin{equation}
\label{ch4:eq:regret-theorem-bound}
\Reg_T(\mathcal M_H)
\le
\frac{\log(H+1)}{\eta}+\eta T\Gamma_\alpha,
\end{equation}
where
\begin{equation}
\label{ch4:eq:def-Gamma-alpha}
\Gamma_\alpha
\coloneqq
\frac{G_\alpha^2}{2}
+
\frac{\alpha G_\alpha(L_x+\alpha L_u)}{(1-\alpha)^2},
\end{equation}
and $G_\alpha=\alpha\bigl(\frac{L_x}{1-\alpha}+L_u\bigr)$ as in
\eqref{ch4:eq:def-surrogate-G}.

If the horizon $T$ and cost bounds $L_x,L_u$ are supplied, compute
$G_\alpha$ and $\Gamma_\alpha$ from the expressions above and choose
$\eta^*\coloneqq\sqrt{\log(H+1)/(T\Gamma_\alpha)}$. Then
\begin{equation}
\label{ch4:eq:regret-theorem-optimized}
\Reg_T(\mathcal M_H)\le 2\sqrt{\Gamma_\alpha T\log(H+1)}.
\end{equation}
\end{Thm:Theorem}

The proof first bounds regret for the surrogate costs, then bounds the
difference between actual and surrogate costs caused by changing the policy
fractions and clipping allocations. This follows the analysis of online
learning with memory in
\cite[App.~C, proof of Thm.~3.1]{kumar2023unbounded}.

\begin{proof}[Proof of Theorem~\ref{ch4:thm:regret-bound}]
Fix $\mathbf m\in\mathcal M_H=S_H$ and use the $\ell_1/\ell_\infty$ dual pair.
Lemma~\ref{ch4:lem:sdac-surrogate} gives convex $\ell_t$ and
$\mathbf g_t\in\partial\ell_t(\mathbf m_t)$ with $\|\mathbf g_t\|_\infty\le G_\alpha$.
By Lemma~\ref{lem:intro:entropy}, $\Psi$ is $1$-strongly convex and
$D_\Psi(\mathbf m,\mathbf m_1)\le\log(H+1)$ at \eqref{ch4:eq:alg-init}.
The updates \eqref{ch4:eq:alg-omd-mult}--\eqref{ch4:eq:alg-omd-proj}
give $\mathbf m_t>\mathbf0$ for $t\in[T+1]$.
Lemma~\ref{lem:intro:mirror-descent} and convexity therefore give
\begin{equation}
\label{ch4:eq:scalar-surrogate-regret}
\begin{aligned}
\sum_{t=1}^T[\ell_t(\mathbf m_t)-\ell_t(\mathbf m)]
&\overset{\mathrm{(A)}}{\le}
 \sum_{t=1}^T\inner{\mathbf g_t}{\mathbf m_t-\mathbf m}\\
&\overset{\mathrm{(B)}}{\le}
 \frac1\eta\sum_{t=1}^T
 [D_\Psi(\mathbf m,\mathbf m_t)-D_\Psi(\mathbf m,\mathbf m_{t+1})]
 +\frac\eta2\sum_{t=1}^T\|\mathbf g_t\|_\infty^2\\
&=\frac{D_\Psi(\mathbf m,\mathbf m_1)
             -D_\Psi(\mathbf m,\mathbf m_{T+1})}{\eta}
       +\frac\eta2\sum_{t=1}^T\|\mathbf g_t\|_\infty^2\\
&\overset{\mathrm{(C)}}{\le}
 \frac{D_\Psi(\mathbf m,\mathbf m_1)}{\eta}
 +\frac\eta2\sum_{t=1}^T\|\mathbf g_t\|_\infty^2\\
&\overset{\mathrm{(D)}}{\le}
 \frac{\log(H+1)}{\eta}+\frac\eta2\sum_{t=1}^TG_\alpha^2\\
&=\frac{\log(H+1)}{\eta}+\frac{\eta TG_\alpha^2}{2}.
\end{aligned}
\end{equation}
Here (A) is the subgradient inequality, (B) sums
\eqref{eq:intro:md-one-step}, (C) uses $D_\Psi\ge0$, and (D) uses
the bounds above. Also, \eqref{eq:intro:md-movement} gives
\begin{equation}
\label{ch4:eq:scalar-md-movement}
\|\mathbf m_{t+1}-\mathbf m_t\|_1
\le\eta\|\mathbf g_t\|_\infty\le\eta G_\alpha,
\qquad t\in[T].
\end{equation}

For any $\mathbf m,\mathbf m'\in\mathcal M_H$, the affine state formula
\eqref{ch4:eq:surr-affine-pair} and H\"older's inequality give
\[
|x_t(\mathbf m)-x_t(\mathbf m')|
=|\inner{\boldsymbol\phi_t^x}{\mathbf m-\mathbf m'}|
\le\|\boldsymbol\phi_t^x\|_\infty\|\mathbf m-\mathbf m'\|_1
\le\frac{\alpha}{1-\alpha}\|\mathbf m-\mathbf m'\|_1.
\]
The final inequality is \eqref{ch4:eq:surr-coefficient-norms}.
Define $v_t\coloneqq u_t(\mathbf m_t)$ and
$e_t\coloneqq|x_t-x_t(\mathbf m_t)|$ for $t\in[T]$. Then $e_1=0$.
Fixed-policy feasibility and \eqref{ch4:eq:alg-feasible-action} give
$0\le v_t\le\alpha x_t(\mathbf m_t)$ and, for $1\le t<T$,
\[
x_{t+1}=[\alpha x_t-v_t]_++w_t,
\qquad
x_{t+1}(\mathbf m_t)=[\alpha x_t(\mathbf m_t)-v_t]_++w_t.
\]
For real $a,b$, $|[a]_+-[b]_+|\le|a-b|$: if both are nonnegative,
the two sides are equal; if both are nonpositive, the left side is zero;
and if they have opposite signs, the left side is the positive argument,
which is at most $|a|+|b|=|a-b|$. Applying this bound gives
\[
\begin{aligned}
|x_{t+1}-x_{t+1}(\mathbf m_t)|
&=|[\alpha x_t-v_t]_+-[\alpha x_t(\mathbf m_t)-v_t]_+|\\
&\le|(\alpha x_t-v_t)-(\alpha x_t(\mathbf m_t)-v_t)|\\
&=\alpha|x_t-x_t(\mathbf m_t)|=\alpha e_t.
\end{aligned}
\]
Hence
\begin{equation}
\label{ch4:eq:scalar-tracking-recursion}
\begin{aligned}
e_{t+1}
&=|x_{t+1}-x_{t+1}(\mathbf m_{t+1})|\\
&\le |x_{t+1}-x_{t+1}(\mathbf m_t)|
     +|x_{t+1}(\mathbf m_t)-x_{t+1}(\mathbf m_{t+1})|\\
&\le\alpha e_t
     +|x_{t+1}(\mathbf m_t)-x_{t+1}(\mathbf m_{t+1})|\\
&\overset{\mathrm{(A)}}{\le}\alpha e_t
    +\frac{\alpha}{1-\alpha}\|\mathbf m_{t+1}-\mathbf m_t\|_1\\
&\overset{\mathrm{(B)}}{\le}\alpha e_t+\frac{\alpha\eta G_\alpha}{1-\alpha}.
\end{aligned}
\end{equation}
Here (A) applies the inequality for
$|x_t(\mathbf m)-x_t(\mathbf m')|$ above at time $t+1$;
(B) uses \eqref{ch4:eq:scalar-md-movement}. Since $e_1=0$,
\[
\begin{aligned}
(1-\alpha)\sum_{t=1}^T e_t+\alpha e_T
&=\left(\sum_{t=1}^T e_t-e_1\right)
 -\alpha\left(\sum_{t=1}^T e_t-e_T\right)\\
&=\sum_{t=1}^{T-1}e_{t+1}-\alpha\sum_{t=1}^{T-1}e_t\\
&=\sum_{t=1}^{T-1}(e_{t+1}-\alpha e_t)\\
&\overset{\mathrm{(C)}}{\le}
 \sum_{t=1}^{T-1}\frac{\alpha\eta G_\alpha}{1-\alpha}
 =\frac{\alpha\eta G_\alpha(T-1)}{1-\alpha}.
\end{aligned}
\]
Here (C) uses \eqref{ch4:eq:scalar-tracking-recursion}. Subtracting
$\alpha e_T$ and dividing by $1-\alpha>0$ gives
\begin{equation}
\label{ch4:eq:scalar-tracking-state-sum}
\sum_{t=1}^T e_t
\le\frac{\alpha\eta G_\alpha(T-1)}{(1-\alpha)^2}
   -\frac{\alpha e_T}{1-\alpha}
\le\frac{\alpha\eta G_\alpha(T-1)}{(1-\alpha)^2}.
\end{equation}
The last inequality uses $e_T\ge0$.
Clipping and fixed-policy feasibility also give
\begin{equation}
\label{ch4:eq:scalar-clipping-error}
0\le v_t-u_t=[v_t-\alpha x_t]_+
\le[\alpha x_t(\mathbf m_t)-\alpha x_t]_+
\le\alpha e_t.
\end{equation}
The first inequality between positive parts uses monotonicity and
$v_t\le\alpha x_t(\mathbf m_t)$. The last uses $[a]_+\le|a|$.
Both state--action pairs lie in $\mathcal D_\alpha$ by feasibility and
\eqref{ch4:eq:model-state-bound}. Apply the cost Lipschitz bound
\eqref{ch4:eq:model-lipschitz} at $(x_t,u_t)$ and
$(x_t(\mathbf m_t),v_t)$, and sum:
\begin{equation}
\label{ch4:eq:scalar-tracking-cost-sum}
\begin{aligned}
\sum_{t=1}^T|c_t(x_t,u_t)-\ell_t(\mathbf m_t)|
&\le\sum_{t=1}^T\bigl(L_xe_t+L_u|u_t-v_t|\bigr)\\
&\le\sum_{t=1}^T\bigl(L_xe_t+\alpha L_ue_t\bigr)\\
&=(L_x+\alpha L_u)\sum_{t=1}^Te_t\\
&\le\frac{\alpha\eta G_\alpha(L_x+\alpha L_u)(T-1)}{(1-\alpha)^2}.
\end{aligned}
\end{equation}
The last two inequalities use \eqref{ch4:eq:scalar-clipping-error}
and \eqref{ch4:eq:scalar-tracking-state-sum}, respectively.

Insert \eqref{ch4:eq:scalar-surrogate-regret} and
\eqref{ch4:eq:scalar-tracking-cost-sum} into the decomposition
\eqref{ch4:eq:regret-decomposition}. Since $T-1\le T$ and
$\sum_{t=1}^T\ell_t(\mathbf m)=J_T(\mathbf m)$, this gives
\[
\begin{aligned}
J_T^{\mathrm{on}}-J_T(\mathbf m)
&\le\frac{\log(H+1)}{\eta}+\frac{\eta TG_\alpha^2}{2}
 +\frac{\alpha\eta G_\alpha(L_x+\alpha L_u)(T-1)}{(1-\alpha)^2}\\
&\le\frac{\log(H+1)}{\eta}+\frac{\eta TG_\alpha^2}{2}
 +\frac{\alpha\eta G_\alpha(L_x+\alpha L_u)T}{(1-\alpha)^2}\\
&=\frac{\log(H+1)}{\eta}
 +\eta T\left[\frac{G_\alpha^2}{2}
       +\frac{\alpha G_\alpha(L_x+\alpha L_u)}{(1-\alpha)^2}\right]\\
&=\frac{\log(H+1)}{\eta}+\eta T\Gamma_\alpha.
\end{aligned}
\]
Taking the supremum over $\mathbf m\in\mathcal M_H$ gives
\eqref{ch4:eq:regret-theorem-bound} by \eqref{ch4:eq:def-regret}.
Since $\log(H+1)>0$ and $\Gamma_\alpha>0$,
$\eta^*=\sqrt{\log(H+1)/(T\Gamma_\alpha)}>0$, and
\[
\frac{\log(H+1)}{\eta^*}
=\eta^*T\Gamma_\alpha
=\sqrt{\Gamma_\alpha T\log(H+1)},
\]
which gives \eqref{ch4:eq:regret-theorem-optimized}.
\end{proof}

\begin{Thm:Remark}[Dependence on the retention coefficient]
\label{ch4:rem:retention-dependence}
Fix $L_x,L_u>0$ and let $\alpha\uparrow1$ through $(0,1)$.
Equation~\eqref{ch4:eq:def-surrogate-G} gives
\[
(1-\alpha)G_\alpha
=\alpha\bigl(L_x+(1-\alpha)L_u\bigr)\longrightarrow L_x.
\]
Multiplying \eqref{ch4:eq:def-Gamma-alpha} by $(1-\alpha)^3$ yields
\[
\begin{aligned}
(1-\alpha)^3\Gamma_\alpha
&=\frac{1-\alpha}{2}\bigl((1-\alpha)G_\alpha\bigr)^2
  +\alpha\bigl((1-\alpha)G_\alpha\bigr)(L_x+\alpha L_u)\\
&\longrightarrow L_x(L_x+L_u)>0.
\end{aligned}
\]
Consequently, the coefficient $2\sqrt{\Gamma_\alpha}$ in the upper bound
\eqref{ch4:eq:regret-theorem-optimized} is
$\Theta((1-\alpha)^{-3/2})$ as $\alpha\uparrow1$.
\end{Thm:Remark}

\subsection{Infinite-memory policies and fixed-fraction comparisons}
\label{ch4:subsec:infinite-memory-sdac}

The infinite-memory policy class is
\begin{equation}
\label{ch4:eq:def-infinite-simplex}
\mathcal M_\infty
\coloneqq\left\{\mathbf m=(m_i)_{i\ge1}:m_i\ge0,\quad
\sum_{i=1}^{\infty}m_i\le1\right\},
\end{equation}
with the feasible action
\begin{equation}
\label{ch4:eq:def-infinite-sdac-policy}
u_t(\mathbf m)\coloneqq\sum_{i=1}^{t-1}m_i\alpha^iw_{t-i}.
\end{equation}
The state follows \eqref{ch4:eq:model-dynamics} from $x_1(\mathbf m)=0$,
and \eqref{ch4:eq:def-fixed-policy-cost} defines $J_T(\mathbf m)$ for
$\mathbf m\in\mathcal M_\infty$ as well.

\begin{Thm:Lemma}[Approximation of infinite-memory SDAC policies]
\label{ch4:lem:sdac-truncation-approx}
Fix positive integers $H,T$ under the model and cost assumptions of
Subsection~\ref{ch4:sec:model-setup}.
For every $\mathbf m\in\mathcal M_\infty$, the infinite-memory policy is
feasible. Its truncation
$\mathbf m^{[H]}\coloneqq(m_1,\ldots,m_H)^\top\in\mathcal M_H$ satisfies
\begin{equation}
\label{ch4:eq:sdac-truncation-cost-gap}
\left|J_T(\mathbf m^{[H]})-J_T(\mathbf m)\right|
\le TG_\alpha\alpha^H.
\end{equation}
If $H\ge T-1$, the two trajectories agree through round $T$, so their
cost difference is zero.
\end{Thm:Lemma}
\begin{proof}
For $t\ge2$, the policy agrees through time $t$ with its truncation at $t-1$,
since $w_{s-i}=0$ for $i\ge s$. The truncated vector belongs to
$\mathcal M_{t-1}$ by \eqref{ch4:eq:def-infinite-simplex} and
\eqref{ch4:eq:def-simplex}, so Lemma~\ref{ch4:lem:sdac-feasibility}
applies with memory $t-1$. At $t=1$, both state and action are zero by
\eqref{ch4:eq:model-init} and \eqref{ch4:eq:def-infinite-sdac-policy}.
This proves feasibility of every $\mathbf m\in\mathcal M_\infty$.

For $t\in[T]$, define
\begin{equation}
\label{ch4:eq:sdac-truncation-gaps}
\begin{aligned}
e_t^u&\coloneqq u_t(\mathbf m)-u_t(\mathbf m^{[H]})
       =\sum_{i=H+1}^{t-1}\alpha^im_iw_{t-i},\\
e_t^x&\coloneqq x_t(\mathbf m^{[H]})-x_t(\mathbf m).
\end{aligned}
\end{equation}
Since $w_s\in[0,1]$, $m_i\ge0$, $\sum_{i\ge1}m_i\le1$, and
$0<\alpha<1$, we have
\[
0\le e_t^u\le\alpha^{H+1}\sum_{i=H+1}^{t-1}m_i
\le\alpha^{H+1}.
\]
Subtracting the two state recursions \eqref{ch4:eq:model-dynamics} gives
\begin{equation}
\label{ch4:eq:sdac-truncation-recursion}
e_{t+1}^x=\alpha e_t^x+e_t^u,\qquad 1\le t<T,\qquad e_1^x=0.
\end{equation}
Iterating and using the action bound yields
\begin{equation}
\label{ch4:eq:sdac-truncation-state-gap}
0\le e_t^x
=\sum_{s=1}^{t-1}\alpha^{t-1-s}e_s^u
\le\alpha^{H+1}\sum_{j=0}^{t-2}\alpha^j
\le\frac{\alpha^{H+1}}{1-\alpha},\qquad t\in[T].
\end{equation}
Both trajectories lie in $\mathcal D_\alpha$
by feasibility and \eqref{ch4:eq:model-state-bound}. The triangle
inequality and \eqref{ch4:eq:model-lipschitz} therefore give
\[
\begin{aligned}
|J_T(\mathbf m^{[H]})-J_T(\mathbf m)|
&\le\sum_{t=1}^T(L_xe_t^x+L_ue_t^u)\\
&\le T\left(\frac{L_x}{1-\alpha}+L_u\right)\alpha^{H+1}
=TG_\alpha\alpha^H.
\end{aligned}
\]
The equality is \eqref{ch4:eq:def-surrogate-G}.
If $H\ge T-1$, the sum defining $e_t^u$ is empty, and the recursion for
$e_t^x$ gives $e_t^u=e_t^x=0$ for all $t\in[T]$.
\end{proof}

For $k\in[0,\alpha]$, the fixed-fraction policy chooses $u_t(k)\coloneqq kx_t(k)$
and starts from $x_1(k)=0$. Substituting this action into
\eqref{ch4:eq:model-dynamics} gives
\begin{equation}
\label{ch4:eq:fixed-fraction-dynamics}
x_{t+1}(k)=(\alpha-k)x_t(k)+w_t,\qquad t\in[T].
\end{equation}

\begin{Thm:Corollary}[Fixed-fraction policies as infinite-memory SDAC policies]
\label{ch4:cor:fraction-as-sdac}
For $k\in[0,\alpha]$, let
\begin{equation}
\label{ch4:eq:fixed-fraction-coefficients}
m_1(k)\coloneqq\frac{k}{\alpha},\qquad
m_{i+1}(k)\coloneqq\left(1-\frac{k}{\alpha}\right)m_i(k),\qquad i\ge1.
\end{equation}
Then $\mathbf m(k)\in\mathcal M_\infty$, and its trajectory
equals that of the feasible fixed-fraction policy $u_t(k)=kx_t(k)$ for every
arrival sequence with entries in $[0,1]$:
\[
\bigl(x_t(\mathbf m(k)),u_t(\mathbf m(k))\bigr)
=\bigl(x_t(k),u_t(k)\bigr),\qquad t\in[T].
\]
\end{Thm:Corollary}
\begin{proof}
Put $q\coloneqq1-k/\alpha\in[0,1]$. The recurrence
\eqref{ch4:eq:fixed-fraction-coefficients} gives
$m_i(k)=(1-q)q^{i-1}\ge0$, with $q^0=1$ also when $q=0$.
For every $N\ge1$,
\[
\sum_{i=1}^N m_i(k)
=\sum_{i=1}^N(q^{i-1}-q^i)
=1-q^N\le1.
\]
Taking $N\to\infty$ shows $\sum_{i\ge1}m_i(k)\le1$.
Thus $\mathbf m(k)\in\mathcal M_\infty$ by
\eqref{ch4:eq:def-infinite-simplex}.
Since $\alpha-k\ge0$ and $w_t\ge0$, the recursion
\eqref{ch4:eq:fixed-fraction-dynamics} from $x_1(k)=0$ gives nonnegative states.
The condition $0\le k\le\alpha$ then gives feasibility in
\eqref{ch4:eq:model-feasibility}. Iterating the state recursion and using
$\alpha^i m_i(k)=k(\alpha-k)^{i-1}$ from
\eqref{ch4:eq:fixed-fraction-coefficients} in
\eqref{ch4:eq:def-infinite-sdac-policy} yields
\[
u_t(\mathbf m(k))
=k\sum_{i=1}^{t-1}(\alpha-k)^{i-1}w_{t-i}
=kx_t(k)=u_t(k).
\]
The two state trajectories start at zero and have identical actions and
arrivals in \eqref{ch4:eq:model-dynamics}, so they also coincide.
\end{proof}

\begin{Thm:Corollary}[No-regret against infinite-memory SDAC policies]
\label{ch4:cor:regret-vs-infinite-sdac}
Under the assumptions of Theorem~\ref{ch4:thm:regret-bound}, run the entropic
SDAC controller with step size $\eta^*=\sqrt{\log(H+1)/(T\Gamma_\alpha)}$
and memory
\begin{equation}
\label{ch4:eq:def-infinite-memory-length}
H=H_T
\coloneqq
\max\!\left\{1,\left\lceil\frac{\log T}{\log(1/\alpha)}\right\rceil\right\}.
\end{equation}
This choice gives $T\alpha^{H_T}\le1$.
For
\begin{equation}
\label{ch4:eq:def-infinite-regret}
\Reg_T(\mathcal M_\infty)
\coloneqq J_T^{\mathrm{on}}-\inf_{\mathbf m\in\mathcal M_\infty}J_T(\mathbf m),
\end{equation}
we have
\begin{equation}
\label{ch4:eq:infinite-sdac-regret-bound}
\Reg_T(\mathcal M_\infty)
\le2\sqrt{\Gamma_\alpha T\log(H_T+1)}+TG_\alpha\alpha^{H_T}.
\end{equation}
For fixed $\alpha,L_x,L_u$, cumulative regret is bounded above by
$O(\sqrt{T\log\log T})$. Consequently, the upper bound on average excess cost
per round tends to zero. This guarantee also applies to the best
fixed-fraction policy chosen in hindsight, because every such policy belongs
to the comparison class $\mathcal M_\infty$:
\[
J_T^{\mathrm{on}}
-\min_{k\in[0,\alpha]}\sum_{t=1}^Tc_t(x_t(k),u_t(k))
\le\Reg_T(\mathcal M_\infty).
\]
\end{Thm:Corollary}
\begin{proof}
For every $\mathbf m\in\mathcal M_\infty$,
$\mathbf m^{[H_T]}\in\mathcal M_{H_T}$ and
\begin{align*}
J_T^{\mathrm{on}}-J_T(\mathbf m)
&=J_T^{\mathrm{on}}-J_T(\mathbf m^{[H_T]})
+J_T(\mathbf m^{[H_T]})-J_T(\mathbf m)\\
&\overset{\mathrm{(A)}}{\le}
 J_T^{\mathrm{on}}-\min_{\mathbf m'\in\mathcal M_{H_T}}J_T(\mathbf m')
 +J_T(\mathbf m^{[H_T]})-J_T(\mathbf m)\\
&=\Reg_T(\mathcal M_{H_T})
 +J_T(\mathbf m^{[H_T]})-J_T(\mathbf m)\\
&\overset{\mathrm{(B)}}{\le}
 2\sqrt{\Gamma_\alpha T\log(H_T+1)}
+J_T(\mathbf m^{[H_T]})-J_T(\mathbf m)\\
&\le2\sqrt{\Gamma_\alpha T\log(H_T+1)}
 +|J_T(\mathbf m^{[H_T]})-J_T(\mathbf m)|\\
&\overset{\mathrm{(C)}}{\le}
 2\sqrt{\Gamma_\alpha T\log(H_T+1)}+TG_\alpha\alpha^{H_T}.
\end{align*}
Here (A) uses $\mathbf m^{[H_T]}\in\mathcal M_{H_T}$, and the next equality
is \eqref{ch4:eq:def-regret}. Step (B) applies
\eqref{ch4:eq:regret-theorem-optimized} with $H=H_T$ and step size $\eta^*$;
(C) is \eqref{ch4:eq:sdac-truncation-cost-gap}. By
\eqref{ch4:eq:def-infinite-regret},
\[
\begin{aligned}
\Reg_T(\mathcal M_\infty)
&=\sup_{\mathbf m\in\mathcal M_\infty}
  [J_T^{\mathrm{on}}-J_T(\mathbf m)]\\
&\le2\sqrt{\Gamma_\alpha T\log(H_T+1)}+TG_\alpha\alpha^{H_T},
\end{aligned}
\]
which is \eqref{ch4:eq:infinite-sdac-regret-bound}. The choice
\eqref{ch4:eq:def-infinite-memory-length} gives
\[
\frac{\log T}{\log(1/\alpha)}\le H_T
\le1+\frac{\log T}{\log(1/\alpha)}.
\]
Since $0<\alpha<1$,
\[
\begin{aligned}
TG_\alpha\alpha^{H_T}
&\le TG_\alpha\alpha^{\log T/\log(1/\alpha)}\\
&=TG_\alpha\exp\!\left(\frac{\log T}{\log(1/\alpha)}\log\alpha\right)\\
&=TG_\alpha e^{-\log T}=G_\alpha.
\end{aligned}
\]
For fixed $\alpha,L_x,L_u$,
$H_T=O(\log T)$ and $\log(H_T+1)=O(\log\log T)$, so the regret bound is
$O(\sqrt{T\log\log T})=o(T)$.
Also,
\begin{align*}
 \inf_{\mathbf m\in\mathcal M_\infty}J_T(\mathbf m)
 &\overset{\mathrm{(A)}}{\le}\min_{k\in[0,\alpha]}J_T(\mathbf m(k))\\
 &=\min_{k\in[0,\alpha]}
   \sum_{t=1}^Tc_t(x_t(k),u_t(k)).
\end{align*}
Here (A) uses $\{\mathbf m(k):k\in[0,\alpha]\}\subseteq\mathcal M_\infty$,
and the equality follows from Corollary~\ref{ch4:cor:fraction-as-sdac}.
Consequently,
\[
\begin{aligned}
J_T^{\mathrm{on}}-\min_{k\in[0,\alpha]}
 \sum_{t=1}^Tc_t(x_t(k),u_t(k))
&\le J_T^{\mathrm{on}}-\inf_{\mathbf m\in\mathcal M_\infty}J_T(\mathbf m)\\
&=\Reg_T(\mathcal M_\infty).
\end{aligned}
\]
\end{proof}

\section{Regret lower bound}
\label{ch4:sec:regret-lower-bound}

We now show that, with $H,\alpha,L_x,L_u$ fixed, no causal feasible controller
can guarantee worst-case expected regret growing more slowly than $\sqrt T$.
Thus the entropic SDAC controller's upper bound has the best possible
dependence on the horizon, up to a constant factor.

\begin{Thm:Theorem}[SDAC regret lower bound]
\label{ch4:thm:sdac-lower-bound}
Fix $\alpha\in(0,1)$, $L_x,L_u>0$, and an integer $H\ge1$.
Assume the dynamics, initial state, and feedback of
Subsection~\ref{ch4:sec:model-setup}, and let the integer horizon satisfy
\begin{equation}
\label{ch4:eq:sdac-lower-horizon}
T\ge2\left\lceil\frac{\log2}{\log(1/\alpha)}\right\rceil+2.
\end{equation}
For every causal feasible online controller $\mathcal A$, possibly
randomized, there is a deterministic sequence of arrivals and costs,
fixed before the interaction, that satisfies the model and cost assumptions
with constants $L_x,L_u$ and gives
\begin{equation}
\label{ch4:eq:sdac-lower-large-horizon}
\mathbb E_{\mathcal A}\!\left[\Reg_T(\mathcal M_H)\right]
\ge\frac{G_\alpha}{8}\sqrt T.
\end{equation}
Here $G_\alpha=\alpha(L_x/(1-\alpha)+L_u)$.
The expectation is over the controller's internal randomization.
The sequence may depend on $\mathcal A$, but not on its realized random
choices.
\end{Thm:Theorem}

The upper bound~\eqref{ch4:eq:regret-theorem-optimized} is
$2\sqrt{\Gamma_\alpha T\log(H+1)}$. For fixed $H,\alpha,L_x,L_u$,
both bounds therefore scale as $\sqrt T$. Their parameter dependences
differ: the upper bound contains $\sqrt{\log(H+1)}$, whereas the lower
bound is independent of $H$. For fixed $L_x,L_u>0$, as $\alpha\uparrow1$
and with $T$ satisfying~\eqref{ch4:eq:sdac-lower-horizon}, the lower bound
is $\Omega((1-\alpha)^{-1}\sqrt T)$, while the upper bound is
$O((1-\alpha)^{-3/2}\sqrt{T\log(H+1)})$ by
Remark~\ref{ch4:rem:retention-dependence}. Thus optimality in $H$ and
$\alpha$ remains unresolved.

\begin{proof}
Fix a causal feasible controller $\mathcal A$.
We use the random-sign argument of Hazan, Agarwal, and Kale
\cite[Sec.~2.3, footnote~3]{hazan2007logarithmic}. See also
\cite[Sec.~3.2, Thm.~3.2]{hazan2016introduction}.

Let $\xi_1,\ldots,\xi_T$ be independent random signs, each equal to $1$
or $-1$ with probability $1/2$, and independent of the controller's
randomization. Write $\boldsymbol\xi\coloneqq(\xi_1,\ldots,\xi_T)$ and set
\[
w_t\coloneqq1,\qquad c_t(x,u)\coloneqq\xi_t(L_xx-L_uu),\qquad t\in[T].
\]
Each cost is affine with gradient $\xi_t(L_x,-L_u)$, so it is convex
and satisfies both the Lipschitz and subgradient bounds
\eqref{ch4:eq:model-lipschitz}--\eqref{ch4:eq:model-subgradient-bounds}.

Since $H\ge1$, the two parameters
\[
\mathbf m^{(0)}\coloneqq\mathbf0,\qquad
\mathbf m^{(1)}\coloneqq\mathbf e_1=(1,0,\ldots,0)^\top
\]
belong to $\mathcal M_H$ and are feasible by
Lemma~\ref{ch4:lem:sdac-feasibility}. Both trajectories start with
state and action zero. For $t=2,\ldots,T$, the policy formula
\eqref{ch4:eq:def-sdac-policy} and dynamics
\eqref{ch4:eq:model-dynamics} give
\[
\begin{aligned}
u_t(\mathbf m^{(0)})&=0,
&x_t(\mathbf m^{(0)})
&=\sum_{s=0}^{t-2}\alpha^s
=\frac{1-\alpha^{t-1}}{1-\alpha},\\
u_t(\mathbf m^{(1)})&=\alpha,
&x_t(\mathbf m^{(1)})&=1.
\end{aligned}
\]
Both policies incur zero cost at $t=1$. Define
\begin{equation}
\label{ch4:eq:sdac-lower-beta}
\beta_t\coloneqq
\alpha\left(L_u+\frac{L_x(1-\alpha^{t-2})}{1-\alpha}\right),
\qquad t=2,\ldots,T.
\end{equation}
Subtracting their costs gives
\[
\begin{aligned}
J_T(\mathbf m^{(0)})-J_T(\mathbf m^{(1)})
&=\sum_{t=2}^T\xi_t
 \left[L_x\left(\frac{1-\alpha^{t-1}}{1-\alpha}-1\right)+\alpha L_u\right]\\
&=\sum_{t=2}^T\xi_t
 \left[L_x\frac{\alpha-\alpha^{t-1}}{1-\alpha}+\alpha L_u\right]\\
&=\sum_{t=2}^T\xi_t\alpha
 \left[L_u+\frac{L_x(1-\alpha^{t-2})}{1-\alpha}\right]\\
&=\sum_{t=2}^T\xi_t\beta_t.
\end{aligned}
\]

By causality, $(x_t,u_t)$ depends only on past signs and the controller's
randomization, and is independent of $\xi_t$. The two fixed-policy
trajectories depend only on the arrivals. Since $\mathbb E[\xi_t]=0$,
\[
\begin{aligned}
\mathbb E_{\boldsymbol\xi,\mathcal A}[J_T^{\mathrm{on}}]
&=\sum_{t=1}^T\mathbb E_{\boldsymbol\xi,\mathcal A}
 [\xi_t(L_xx_t-L_uu_t)]\\
&=\sum_{t=1}^T\mathbb E[\xi_t]\,
 \mathbb E_{\boldsymbol\xi,\mathcal A}[L_xx_t-L_uu_t]=0,\\
\mathbb E_{\boldsymbol\xi}[J_T(\mathbf m^{(j)})]
&=\sum_{t=1}^T\mathbb E[\xi_t]\,
 [L_xx_t(\mathbf m^{(j)})-L_uu_t(\mathbf m^{(j)})]=0,
 \quad j\in\{0,1\}.
\end{aligned}
\]

Using $\min\{a,b\}=(a+b-|a-b|)/2$,
\[
\begin{aligned}
\mathbb E_{\boldsymbol\xi,\mathcal A}[\Reg_T(\mathcal M_H)]
&\overset{\mathrm{(A)}}{\ge}
 \mathbb E_{\boldsymbol\xi,\mathcal A}
 \bigl[J_T^{\mathrm{on}}-\min\{J_T(\mathbf m^{(0)}),J_T(\mathbf m^{(1)})\}\bigr]\\
&=\mathbb E_{\boldsymbol\xi,\mathcal A}[J_T^{\mathrm{on}}]
 -\frac12\mathbb E_{\boldsymbol\xi}
 \bigl[J_T(\mathbf m^{(0)})+J_T(\mathbf m^{(1)})\bigr]\\
&\quad+\frac12\mathbb E_{\boldsymbol\xi}
 \left|J_T(\mathbf m^{(0)})-J_T(\mathbf m^{(1)})\right|\\
&\overset{\mathrm{(B)}}{=}\frac12\mathbb E_{\boldsymbol\xi}
 \left|J_T(\mathbf m^{(0)})-J_T(\mathbf m^{(1)})\right|\\
&=\frac12\mathbb E_{\boldsymbol\xi}
\left|\sum_{t=2}^T\xi_t\beta_t\right|\\
&\overset{\mathrm{(C)}}{\ge}\frac{1}{2\sqrt2}\sqrt{\sum_{t=2}^T\beta_t^2}.
\end{aligned}
\]
Here (A) uses \eqref{ch4:eq:def-regret} and
$\{\mathbf m^{(0)},\mathbf m^{(1)}\}\subseteq\mathcal M_H$;
(B) uses the zero expected costs, and the next equality substitutes
their difference above;
(C) is the $p=1$ Khintchine inequality \cite{szarek1976best} for the
deterministic coefficients $\beta_2,\ldots,\beta_T$ and independent signs.

Let $q_\alpha\coloneqq\lceil\log2/\log(1/\alpha)\rceil$.
For $t\ge q_\alpha+2$, $0<\alpha<1$ gives
\[
\alpha^{t-2}\le\alpha^{q_\alpha}
\le\alpha^{\log2/\log(1/\alpha)}=\frac12.
\]
Thus \eqref{ch4:eq:sdac-lower-beta} gives
\[
\begin{aligned}
\beta_t
&=\alpha\left(L_u+\frac{L_x(1-\alpha^{t-2})}{1-\alpha}\right)\\
&\ge\alpha L_u+\frac{\alpha L_x}{2(1-\alpha)}\\
&\ge\frac{\alpha L_u}{2}+\frac{\alpha L_x}{2(1-\alpha)}
 =\frac{G_\alpha}{2}.
\end{aligned}
\]
By \eqref{ch4:eq:sdac-lower-horizon}, $q_\alpha+1\le T/2$, so
\[
\begin{aligned}
\sum_{t=2}^T\beta_t^2
&\ge\sum_{t=q_\alpha+2}^T\beta_t^2\\
&\ge(T-q_\alpha-1)\left(\frac{G_\alpha}{2}\right)^2\\
&\ge\frac T2\left(\frac{G_\alpha}{2}\right)^2.
\end{aligned}
\]
Consequently,
\[
\frac{1}{2\sqrt2}\sqrt{\sum_{t=2}^T\beta_t^2}
\ge\frac{1}{2\sqrt2}\sqrt{\frac T2\left(\frac{G_\alpha}{2}\right)^2}
=\frac{G_\alpha}{8}\sqrt T.
\]
The expectation over $\boldsymbol\xi$ averages over $2^T$ equally likely
sign sequences. At least one sequence attains or exceeds this average.
Fixing it before the interaction gives
\begin{equation}
\label{ch4:eq:sdac-lower-exact}
\mathbb E_{\mathcal A}\!\left[\Reg_T(\mathcal M_H)\right]
\ge\frac{1}{2\sqrt2}\sqrt{\sum_{t=2}^T\beta_t^2}
\ge\frac{G_\alpha}{8}\sqrt T,
\end{equation}
Equation~\eqref{ch4:eq:sdac-lower-exact} proves
\eqref{ch4:eq:sdac-lower-large-horizon}.
\end{proof}

\section{Conclusion}
\label{ch4:sec:conclusion}

Theorem~\ref{ch4:thm:regret-bound} bounds the regret of the feasible entropic
SDAC controller by $2\sqrt{\Gamma_\alpha T\log(H+1)}$ with the optimized
step size. With $H=H_T$ from
\eqref{ch4:eq:def-infinite-memory-length},
Corollary~\ref{ch4:cor:regret-vs-infinite-sdac} gives
$O(\sqrt{T\log\log T})$ regret against infinite-memory SDAC policies
for fixed $\alpha,L_x,L_u$, including every feasible scalar fixed-fraction
policy.

Theorem~\ref{ch4:thm:sdac-lower-bound} shows that, with $H,\alpha$ and
$L_x,L_u>0$ fixed, no causal feasible controller can guarantee worst-case
expected regret growing more slowly than $\sqrt T$. The upper bound therefore
has the best possible dependence on the horizon, up to a constant factor.
It contains a factor $\sqrt{\log(H+1)}$ that is absent from the lower bound.
As $\alpha$ approaches one, the bounds also differ by a factor
$(1-\alpha)^{-1/2}$, subject to the horizon condition
\eqref{ch4:eq:sdac-lower-horizon}. We do not yet know whether these additional
factors in the upper bound are necessary.

The upper bounds require revelation of the full cost function and enough
capacity to avoid saturation, as specified in
Subsection~\ref{ch4:sec:model-setup}.
The learning update requires a cost subgradient at the state and action that
would result from using the current policy fractions from the start. These
can differ from the actual storage level and allocation. Feedback only at
the actual state and action may therefore be insufficient to compute the
update. If storage capacity causes overflow, the storage level may cease to
be affine in the policy fractions, so the convexity argument in
Lemma~\ref{ch4:lem:sdac-surrogate} no longer applies.

\section{Action clipping as a projection}
\label{ch4:subsec:scalar-entropy}

\begin{Thm:Lemma}[Action clipping as a Bregman projection]
\label{ch4:lem:availability-projection}
For $t\in[T]$, define the set of parameters compatible with the current state by
\begin{equation}
\label{ch4:eq:availability-set}
\mathcal{B}_t
\coloneqq
\left\{
\mathbf{m}\in\mathcal{M}_H:
\inner{\boldsymbol{\phi}_t^u}{\mathbf{m}}
\le \alpha x_t
\right\}.
\end{equation}
Let $\mathbf{m}_t^{\mathrm{exec}}$ be the Bregman projection of
$\mathbf{m}_t$ onto $\mathcal{B}_t$:
\begin{equation}
\label{ch4:eq:availability-projection}
\mathbf{m}_t^{\mathrm{exec}}
\in
\argmin_{\mathbf{m}\in\mathcal{B}_t}
D_\Psi(\mathbf{m},\mathbf{m}_t).
\end{equation}
Then
\[
\inner{\boldsymbol{\phi}_t^u}{\mathbf{m}_t^{\mathrm{exec}}}
=
\min\!\left\{
\inner{\boldsymbol{\phi}_t^u}{\mathbf{m}_t},
\alpha x_t
\right\}
=u_t.
\]
\end{Thm:Lemma}

\begin{proof}
Write $a_t\coloneqq\inner{\boldsymbol\phi_t^u}{\mathbf m_t}$ and
$b_t\coloneqq\inner{\boldsymbol\phi_t^u}{\mathbf m_t^{\mathrm{exec}}}$.
If $a_t\le\alpha x_t$, then $\mathbf m_t\in\mathcal B_t$. In this case,
\eqref{eq:intro:entropy-curvature} and
$D_\Psi(\mathbf m_t,\mathbf m_t)=0$ give
$\mathbf m_t^{\mathrm{exec}}=\mathbf m_t$.
Otherwise $a_t>\alpha x_t\ge b_t$. Suppose $b_t<\alpha x_t$ and choose
$0<\theta<(\alpha x_t-b_t)/(a_t-b_t)<1$.
The point
$\mathbf m_\theta\coloneqq(1-\theta)\mathbf m_t^{\mathrm{exec}}+\theta\mathbf m_t$
belongs to $\mathcal M_H$ by convexity, and
$\inner{\boldsymbol\phi_t^u}{\mathbf m_\theta}
=b_t+\theta(a_t-b_t)<\alpha x_t$, so $\mathbf m_\theta\in\mathcal B_t$.
Convexity of $D_\Psi(\cdot,\mathbf m_t)$ gives
\[
\begin{aligned}
D_\Psi(\mathbf m_\theta,\mathbf m_t)
&\le(1-\theta)D_\Psi(\mathbf m_t^{\mathrm{exec}},\mathbf m_t)
   +\theta D_\Psi(\mathbf m_t,\mathbf m_t)\\
&=(1-\theta)D_\Psi(\mathbf m_t^{\mathrm{exec}},\mathbf m_t)
<D_\Psi(\mathbf m_t^{\mathrm{exec}},\mathbf m_t).
\end{aligned}
\]
The strict inequality uses $\theta>0$ and \eqref{eq:intro:entropy-curvature}, since
$a_t\ne b_t$. This contradicts optimality, so $b_t=\alpha x_t$, as required.
\end{proof}
